\documentclass[12pt]{article}
\usepackage[utf8]{inputenc}
\usepackage[T1]{fontenc}
\usepackage[margin=1in]{geometry}
\usepackage{setspace}
\usepackage{xcolor}
\usepackage[colorlinks=true, urlcolor=blue!60!black, linkcolor=blue!60!black]{hyperref}
\setlength{\parindent}{2em}
\setlength{\parskip}{1.25em}
\renewcommand{\baselinestretch}{1.0}
\begin{document}
\begin{center}
{\LARGE \textit{McDonald v. City of Chicago}, 561 U.S. 742 (2010)\par}
\vspace{0.8em}
{\large Case Brief --- Four-Page Edition\par}
\end{center}
\vspace{1em}

\section{Memory Jogger}

In a 5--4 decision, the Court held that the Fourteenth Amendment incorporates the Second Amendment right to keep and bear arms for self-defense against state and local governments, with a four-Justice plurality relying on the Due Process Clause and Justice Thomas concurring on the Privileges or Immunities Clause (\href{https://www.courtlistener.com/opinion/149702/mcdonald-v-city-of-chicago/}{561 U.S. 742}).

\section{Detailed Case Facts}

Two years after \textit{District of Columbia v. Heller} recognized an individual right to keep a handgun in the home for self-defense against federal infringement, four Chicago residents---Otis McDonald, Adam Orlov, and Colleen and David Lawson---challenged handgun bans in Chicago and neighboring Oak Park. Chicago barred possessing any firearm without a registration certificate while refusing to register almost all handguns, effectively banning them; Oak Park made private possession of any firearm unlawful. Chicago had adopted its ban in 1982 to protect residents from firearm injury and death. McDonald, a community activist in his late seventies, had been threatened by drug dealers, and Colleen Lawson's home had been targeted by burglars; the petitioners kept handguns outside city limits but could not keep them at home. Because \textit{Heller} involved a federal enclave, the case turned on whether the Second Amendment binds the States.

\section{Procedural History}

The petitioners, joined by gun-rights groups, sued in the Northern District of Illinois alongside parallel NRA suits. The district court dismissed, citing Seventh Circuit precedent upholding a municipal handgun ban and \textit{Heller}'s reservation of incorporation. The Seventh Circuit (Easterbrook, C.J.) affirmed: though it called the reasoning of \textit{Cruikshank}, \textit{Presser}, and \textit{Miller v. Texas} ``defunct,'' only the Supreme Court could overrule them. The Supreme Court granted certiorari and decided the case on June 28, 2010.

\section{Judicial Votes}

Reversed and remanded, \textbf{5--4}.
\begin{itemize}
\item \textbf{Opinion of the Court (in part):} Justice Alito, joined by Chief Justice Roberts and Justices Scalia, Kennedy, and Thomas.
\item \textbf{Plurality (Parts II-C, IV, V):} Alito, joined by Roberts, Scalia, and Kennedy (Due Process incorporation; declining to revisit \textit{Slaughter-House}).
\item \textbf{Concurrences:} Justice Scalia; Justice Thomas, concurring in part and in the judgment (Privileges or Immunities).
\item \textbf{Dissents:} Justice Stevens; Justice Breyer, joined by Justices Ginsburg and Sotomayor.
\end{itemize}

\section{Holding}

The Fourteenth Amendment makes the Second Amendment right recognized in \textit{Heller} fully applicable to the States. A Bill of Rights guarantee is incorporated if it is ``fundamental to our scheme of ordered liberty'' or ``deeply rooted in this Nation's history and tradition,'' and it then applies to the States as it does to the federal government. No five Justices agreed on which clause accomplishes incorporation. Part IV's repetition of \textit{Heller}'s assurances about felon bans, sensitive places, and commercial sales, and the footnote listing unincorporated rights, are dicta.

\section{Analysis of Opinions}

\textbf{Alito (Court and plurality).} Tracing incorporation from \textit{Barron} to selective incorporation, the majority found self-defense ``the central component'' of the right and deeply rooted from the English Bill of Rights through Reconstruction, stressing Southern efforts to disarm freedmen, the Freedmen's Bureau Act's protection of ``the constitutional right to bear arms,'' and the 22 of 37 States protecting the right in 1868. It rejected a mere antidiscrimination reading of Section 1. The plurality declined to revive Privileges or Immunities, and rejected a ``civilized'' systems test, a watered-down state standard, and interest-balancing.

\textbf{Scalia (concurring).} Though skeptical of substantive due process, he accepted incorporation as ``both long established and narrowly limited'' and attacked Stevens's approach as an unconstrained ``living Constitution'' with ``omnidirectional guideposts.''

\textbf{Thomas (concurring in part and in the judgment).} Substantive due process is a ``legal fiction''; the 1868 public understood the privileges or immunities of national citizenship to include the right to keep and bear arms, so \textit{Cruikshank} and \textit{Slaughter-House} should be rejected.

\textbf{Stevens (dissenting).} The liberty clause need not mirror the Bill of Rights ``jot-for-jot''; handgun ownership is not implicit in ordered liberty, given firearms' ``fundamentally ambivalent relationship to liberty'' and the value of local experimentation.

\textbf{Breyer (dissenting).} ``Nothing in the Second Amendment's text, history, or underlying rationale'' makes private armed self-defense fundamental; lack of consensus, no political-process failure, federalism, and judicial incompetence on gun policy counsel against incorporation.

\section{Examples: Future Applications}

\begin{enumerate}
\item \textbf{Same-Side (Invalid): Public-Housing Handgun Ban.} A city ordinance barring tenants from keeping any operable handgun at home is unconstitutional because the right applies fully to local governments.
\item \textbf{Same-Side (Incorporated): Excessive Fines.} Seizing a \$42,000 vehicle for an offense with a \$10,000 maximum fine implicates a right incorporated under the same test, as \textit{Timbs v. Indiana} later held.
\item \textbf{Opposite-Side (Valid): Felon Possession Ban.} \textit{McDonald} repeated \textit{Heller}'s assurance that longstanding bans on possession by felons remain valid.
\item \textbf{Opposite-Side (Not Incorporated): Grand Jury.} A State may prosecute felonies by information; the plurality left \textit{Hurtado} undisturbed.
\item \textbf{Fence-Sitter: Lawful Permanent Residents.} A ban on handguns for non-citizen residents exposes the 4--1--4 split: due process protects ``any person,'' but Thomas's clause protects only ``citizens of the United States.''
\end{enumerate}

\section{Critique}

\textbf{Originalist views.} Originalists praised the result but criticized the refusal to overrule \textit{Slaughter-House}, which scholars across the spectrum think misread the Privileges or Immunities Clause; they still debate whether reviving it would unleash unenumerated rights and whether 1791 or 1868 meaning controls.

\textbf{Progressive and doctrinal views.} Critics say \textit{McDonald} entrenched \textit{Heller}'s contested ``law-office history,'' stripped cities of tools against gun violence, and gave no standard of scrutiny. Scholars also note the ``voting paradox'': five Justices rejected Due Process incorporation and eight rejected Privileges or Immunities incorporation.

\section{Key Quotations}

\begin{quote}
``In \textit{Heller}, we held that the Second Amendment protects the right to possess a handgun in the home for the purpose of self-defense. Unless considerations of \textit{stare decisis} counsel otherwise, a provision of the Bill of Rights that protects a right that is fundamental from an American perspective applies equally to the Federal Government and the States. See \textit{Duncan}, 391 U.S., at 149, and n.~14. We therefore hold that the Due Process Clause of the Fourteenth Amendment incorporates the Second Amendment right recognized in \textit{Heller}.''
\end{quote}
\noindent --- Justice Alito, Plurality Opinion, 561 U.S. at 791.

\begin{quote}
``The notion that a constitutional provision that guarantees only `process' before a person is deprived of life, liberty, or property could define the substance of those rights strains credulity for even the most casual user of words.''
\end{quote}
\noindent --- Justice Thomas, concurring in part and in the judgment, 561 U.S. at 811.

\begin{quote}
``Self-defense is a basic right, recognized by many legal systems from ancient times to the present day, and in \textit{Heller}, we held that individual self-defense is `the central component' of the Second Amendment right.''
\end{quote}
\noindent --- Justice Alito, Opinion of the Court, 561 U.S. at 767.

\end{document}
